\documentclass[journal]{IEEEtran}
\usepackage{cite}
\usepackage{amsmath,amssymb}
\usepackage{graphicx}
\usepackage{booktabs}
\usepackage{url}

\begin{document}

\title{Explainable Machine Learning for Predicting Longitudinal Motor Function Outcomes in Spinal Muscular Atrophy Type 1 and Type 2 Using Baseline Clinical, Genetic, and CHOP-INTEND/HFMSE Features}

\author{Aarav~Sharma,~\IEEEmembership{Member,~IEEE,}
        Elena~Rostova,
        Marcus~Vance,
        Karis~Patel,
        Sofia~Lindqvist,
        and~Julian~Thorne%
\thanks{A. Sharma, K. Patel, and J. Thorne are with the Department of Pediatric Neurology and Neuromuscular Medicine, University Medical Center, Zurich, Switzerland (e-mail: a.sharma@neurology-med.org).}%
\thanks{E. Rostova is with the Institute of Medical Artificial Intelligence and Biomedical Engineering, Technical University of Munich, Munich, Germany.}%
\thanks{M. Vance is with the Department of Neurology and Neurobiology, Johns Hopkins University School of Medicine, Baltimore, MD, USA.}%
\thanks{S. Lindqvist is with the Division of Clinical Genetics, Karolinska University Hospital, Stockholm, Sweden.}%
\thanks{Manuscript received March 15, 2026; revised May 20, 2026; accepted June 10, 2026.}}

\markboth{IEEE Transactions on Biomedical Engineering,~Vol.~73, No.~8, August~2026}%
{Sharma \MakeLowercase{\textit{et al.}}: Explainable ML for Longitudinal Motor Function in SMA}

\maketitle

\begin{abstract}
Spinal Muscular Atrophy (SMA) is an autosomal recessive neurodegenerative disorder caused by biallelic loss of \textit{SMN1}, leading to progressive spinal $\alpha$-motor neuron loss, muscular atrophy, and quadriparesis. Although revolutionary disease-modifying therapies (DMTs)---nusinersen, risdiplam, and onasemnogene abeparvovec---have redefined the natural history, longitudinal motor responses remain markedly heterogeneous. Clinicians lack rigorous, interpretable prognostic instruments to forecast multi-year motor trajectory outcomes and guide tailored clinical decisions. Here, we present a multicenter explainable machine learning (ML) framework for predicting 24-month longitudinal motor score changes ($\Delta\text{Score}_{24\text{m}}$) in 350 pediatric SMA patients (SMA Type 1: $n=154$, evaluated via CHOP-INTEND; SMA Type 2: $n=196$, evaluated via HFMSE). Using 11 baseline clinical, genetic (\textit{SMN2} copies), electrophysiological (compound muscle action potential, CMAP), anthropometric, and psychometric features, six machine learning architectures were trained within a 10-fold repeated nested cross-validation paradigm. The Gradient Boosting Regressor (GBM) exhibited superior prognostic precision, achieving a cross-validated $R^2$ of 0.840, an RMSE of 1.77 points, an MAE of 1.37 points, and a Pearson correlation $r = 0.926$ ($p < 0.001$). Bland-Altman agreement confirmed zero systematic bias (mean difference: $+0.00$ points; 95\% limits of agreement: $-3.46$ to $+3.46$). Game-theoretic Shapley Additive exPlanations (SHAP TreeExplainer) revealed that therapeutic initiation delay is the single most dominant determinant of 24-month recovery (mean $|\text{SHAP}| = 2.14$ points), followed by baseline CMAP amplitude (1.82 points), baseline functional score (1.45 points), and \textit{SMN2} copy number (1.28 points). Partial dependence functions identified steep non-linear therapeutic windows, wherein every month of delayed intervention beyond 3.5 months in Type 1 and 8.0 months in Type 2 precipitously attenuated prospective motor gains. Individualized waterfall decompositions demonstrated bedside utility by explaining divergent therapeutic trajectories. This framework establishes an accurate, transparent, and clinically actionable paradigm for precision prognosis and trial stratification in pediatric neurology.
\end{abstract}

\begin{IEEEkeywords}
Spinal Muscular Atrophy, Explainable Artificial Intelligence, SHAP, CHOP-INTEND, HFMSE, Longitudinal Outcomes, Precision Neurology, Disease-Modifying Therapy, Prognostic Machine Learning.
\end{IEEEkeywords}

\section{Introduction}

\IEEEPARstart{S}{pinal} Muscular Atrophy (SMA) is an autosomal recessive, debilitating neuromuscular disorder characterized by the progressive degeneration of lower motor neurons within the anterior horns of the spinal cord and lower brainstem \cite{mercuri2018nusinersen,finkel2018diagnosis}. Affecting approximately 1 in 6,000 to 1 in 10,000 live births globally, SMA represents the leading genetic etiology of infant mortality and lifelong pediatric disability \cite{prior2004homozygous}. At the molecular level, more than 95\% of all cases are attributable to homozygous deletions, gene conversions, or loss-of-function point mutations in the \textit{Survival Motor Neuron 1} (\textit{SMN1}) gene located on chromosome 5q13.2 \cite{lefebvre1995identification}. Deficient expression of the full-length SMN protein disrupts the assembly of small nuclear ribonucleoprotein (snRNP) complexes, alters fundamental pre-mRNA alternative splicing machinery, and compromises axonal mRNA transport along the motor axon terminal arborization \cite{pellizzoni1998novel,wirth2021spinal}.

\subsection{Genomic Modifiers and Alternative Splicing}
The clinical expression of 5q-SMA is uniquely modified by a paralogous centromeric gene, \textit{SMN2}, which arose through an evolutionary inverted duplication on chromosome 5q. Although nearly identical in sequence, \textit{SMN2} harbors a crucial silent transition (c.840C$>$T) within exon 7 \cite{lorson1999single}. This single base change disrupts an exonic splicing enhancer recognized by ASF/SF2 while simultaneously generating an exonic splicing silencer recognized by heterogeneous nuclear ribonucleoprotein A1 (hnRNP A1). Consequently, approximately 85\% to 90\% of pre-mRNA transcripts transcribed from \textit{SMN2} omit exon 7, encoding a truncated, structurally unstable protein isoform ($\text{SMN}\Delta7$) that undergoes rapid degradation via the ubiquitin-proteasome degradation cascade. Only 10\% to 15\% of transcripts from \textit{SMN2} yield functional, full-length SMN protein. As a result, the genomic copy number of \textit{SMN2} serves as the primary, yet imperfect, biological modifier of SMA phenotypic severity.

\subsection{Molecular Architecture of the SMN Complex}
The survival motor neuron protein functions not as an isolated monomer, but as the central nucleating scaffold of a macromolecular multiprotein assembly termed the SMN complex. Comprising SMN, Gemins 2 through 8, and Unrip, this oligomeric machine is essential for the chaperoned assembly of heptameric Sm-protein rings onto uridylate-rich small nuclear RNAs (snRNAs U1, U2, U4, U5) to generate mature spliceosomal snRNPs \cite{pellizzoni1998novel}. In the selective vulnerability of motor neurons, SMN depletion impairs snRNP biogenesis, disproportionately triggering alternative splicing errors across genes with minor (U12-dependent) introns, including genes governing mitochondrial energetics, calcium homeostasis, and synaptic transmission \cite{wirth2021spinal}. Furthermore, SMN plays an extra-spliceosomal role in motor axons, forming ribonucleoprotein granules that transport mRNAs encoding actin-regulatory proteins (such as $\beta$-actin and $\beta$-tropomyosin) to distal growth cones and terminal presynaptic boutons. Loss of axonal SMN precipitates premature growth cone collapse, terminal arborization arrest, and failure of neuromuscular junction (NMJ) maturation.

\subsection{Phenotypic Taxonomy: SMA Type 1 vs. Type 2}
Historically, SMA has been classified into discrete clinical subtypes determined by the age of onset and the highest achieved gross motor milestone:
\begin{enumerate}
    \item \textbf{SMA Type 1 (Werdnig-Hoffmann Disease):} Represents approximately 60\% of incident cases, with symptom onset occurring prior to 6 months of age (median $\sim$2.5 months). Infants present with generalized severe hypotonia (the classical ``floppy infant'' appearance), symmetrical flaccid quadriparesis, bell-shaped thoracic deformity secondary to diaphragmatic breathing with intercostal paralysis, absent tendon reflexes, and fasciculations of the tongue \cite{devivo2021clinical}. Without disease-modifying treatment, natural history studies document that affected infants never achieve unsupported sitting, and over 90\% succumb to respiratory failure or require invasive mechanical ventilation before their second birthday \cite{sproule2009increased}.
    \item \textbf{SMA Type 2 (Dubowitz Disease):} Represents roughly 25\% of cases, with clinical manifestations emerging between 6 and 18 months of age. These children acquire the capacity to sit upright unsupported when placed, but fail to achieve independent upright ambulation \cite{mercuri2018diagnosis}. While survival typically extends into adulthood with proactive pulmonary and orthopedic support, patients face substantial morbidity, including accelerated neuromuscular scoliosis, joint contractures, restrictive ventilatory insufficiency, and marked proximal limb weakness \cite{coratti2021longitudinal}.
\end{enumerate}

\subsection{Psychometric Motor Scales: CHOP-INTEND and HFMSE}
To track disease progression and therapeutic efficacy across distinct developmental stages, specialized, validated motor assessment batteries are routinely deployed:
\begin{enumerate}
    \item \textbf{CHOP-INTEND (Children's Hospital of Philadelphia Infant Test of Neuromuscular Disorders):} Standardized for non-sitting infants with profound weakness, evaluating 14 active, antigravity, and reflexive motor behaviors (total range: 0--64 points) \cite{glanzman2010chop}. In clinical trials and real-world registries, a gain of $\ge +4$ points is widely established as the Minimal Clinically Important Difference (MCID) \cite{finkel2021endear,connolly2016clinical}.
    \item \textbf{HFMSE (Hammersmith Functional Motor Scale Expanded):} Tailored for sitters and ambulatory individuals, consisting of 33 discrete motor tasks (each scored 0--2; total range: 0--66 points) evaluating rolling, prop-sitting, four-point kneeling, crawling, and transitional maneuvers \cite{ohagen2007expanded}. An improvement of $\ge +3$ points denotes the validated MCID \cite{montes2019minimal}.
\end{enumerate}

\subsection{The Therapeutic Revolution and Clinical Prognostic Dilemma}
Over the past decade, three revolutionary Disease-Modifying Therapies (DMTs) have received international regulatory authorization:
\begin{itemize}
    \item \textbf{Nusinersen (Spinraza):} An intrathecally administered 2'-O-methoxyethyl antisense oligonucleotide that targets the intronic splicing silencer N1 (ISS-N1) within \textit{SMN2} intron 7, displacing hnRNP A1/A2 and driving inclusion of exon 7 \cite{hua2007enhancement,mercuri2018nusinersen}.
    \item \textbf{Risdiplam (Evrysdi):} An orally bioavailable, systemically distributed small-molecule splicing modulator that selectively shifts the 5' splice site of \textit{SMN2} exon 7 toward inclusion, elevating functional SMN protein across both central nervous tissue and peripheral organs \cite{ratni2018discovery,mercuri2021risdiplam,chiriboga2023sunfish}.
    \item \textbf{Onasemnogene Abeparvovec (Zolgensma):} An adeno-associated virus serotype 9 (AAV9) gene replacement vector delivering an intact human \textit{SMN1} cDNA under a chicken-$\beta$-actin hybrid promoter to motor neurons via a single intravenous infusion \cite{mendell2017single,strauss2022onasemnogene,servais2022onasemnogene}.
\end{itemize}

While these therapies have converted previously fatal infantile SMA into a survivable condition with milestone gains, real-world registry evidence from 2020--2026 demonstrates striking heterogeneity in clinical trajectory. While some infants rapidly gain head control, sitting, and independent ambulation, a substantial proportion display modest functional gains, persistent respiratory dependency, or premature plateauing \cite{finkel2022restore,hagenacker2020nusinersen,pechmann2023effect,wurster2023long}. Predicting the individual 24-month trajectory remains an unresolved challenge for pediatric neurologists.

\subsection{The Human and Socioeconomic Burden of Delay}
In chronic progressive neuromuscular diseases, delays in diagnosis and therapeutic initiation inflict immense personal, psychological, and financial costs on affected families and healthcare ecosystems. Health-economic evaluations across European and North American healthcare systems indicate that patients requiring lifelong mechanical ventilation and intensive supportive care generate direct lifetime medical costs exceeding \$4.5 million. Furthermore, irreversible loss of anterior horn motor neurons during diagnostic Odysseys permanently truncates developmental milestones, placing overwhelming caregiver burdens on parents. Providing objective algorithmic tools to quantify the exact cost of therapeutic delay provides compelling justification for universal newborn screening and immediate DMT access.

\subsection{Research Gaps and Objectives}
Existing prognostic models in SMA suffer from critical shortcomings:
\begin{enumerate}
    \item \textbf{Univariate Oversimplification:} Clinical prognosis is frequently based almost entirely on \textit{SMN2} copy number or chronological age at presentation, ignoring non-linear interactions with baseline electrophysiology (CMAP), nutritional state, and existing functional motor reserve \cite{devivo2024compound,tizzano2023biomarkers}.
    \item \textbf{Short-Term Horizons:} Most published computational investigations restrict predictions to 6-month horizons, failing to anticipate the 24-month window where therapeutic deceleration and plateaus manifest.
    \item \textbf{The Opaque ``Black-Box'' Barrier:} Conventional deep neural networks or complex ensembles lack transparency, preventing clinicians from understanding the biological rationale behind algorithmic predictions \cite{rudin2019stop,lundberg2017unified}.
    \item \textbf{Scale Incompatibility:} Previous works rarely integrate disparate pediatric motor instruments (CHOP-INTEND vs. HFMSE) into a mathematically unified prognostic architecture.
\end{enumerate}

To bridge these gaps, this investigation formulates and validates an Explainable Machine Learning framework for forecasting 24-month longitudinal motor score improvements ($\Delta\text{Score}_{24\text{m}}$) in SMA Type 1 and Type 2. We leverage game-theoretic Shapley Additive exPlanations (SHAP) to unpack global feature hierarchies, delineate critical non-linear therapeutic windows, and generate patient-specific bedside explanations.

\begin{figure*}[!t]
\centering
\includegraphics[width=\textwidth]{SMA_ML_Longitudinal_Figures/Figure_01_Clinical_and_Pathophysiological_Overview.png}
\caption{Comprehensive Clinical and Pathophysiological Overview of Spinal Muscular Atrophy Type 1 and Type 2. (A) Chromosomal 5q13.2 genomic architecture illustrating \textit{SMN1} vs. \textit{SMN2} alternative splicing governed by the c.840C$>$T transition. (B) Neuropathological sequence from nuclear snRNP depletion to spinal anterior horn $\alpha$-motor neuron apoptosis, denervation, and muscular atrophy. (C) Clinical taxonomy comparing Type 1 (Werdnig-Hoffmann) vs. Type 2 (Dubowitz) across onset, milestones, and natural history. (D) Molecular mechanisms of disease-modifying therapies (Nusinersen, Risdiplam, Onasemnogene Abeparvovec).}
\label{fig:overview}
\end{figure*}

\section{Related Work and Literature Review (2020--2026)}

\subsection{Real-World Registries and Longitudinal Trajectories}
The introduction of DMTs sparked the establishment of major international observational registries. The multinational RESTORE registry \cite{finkel2022restore} and the European SMArtCARE network \cite{hagenacker2020nusinersen,pechmann2023effect,wurster2023long} have documented multi-year motor outcomes across thousands of patients. Coratti \textit{et al.} (2021, 2023) demonstrated that although motor gains on the HFMSE can be maintained over several years in Type 2 and 3 patients, the annualized velocity of improvement diminishes significantly after the initial 12 months \cite{coratti2021longitudinal,coratti2023predictors}. In infant cohorts evaluated using CHOP-INTEND, Pane \textit{et al.} (2022) reported that baseline motor score and age at therapeutic onset represent pivotal determinants of milestone attainment \cite{pane2022long}. However, registry publications have primarily relied on standard Kaplan-Meier estimates and multivariable linear regressions, which fail to capture high-order non-linearities and multivariate interactions.

\subsection{Electrophysiological and Molecular Biomarkers}
While \textit{SMN2} copy number remains the paramount genetic prognosticator, clinical literature from 2021 to 2025 emphasizes that copy count alone is an incomplete predictor of individual motor recovery \cite{wirth2020twenty}. De Vivo \textit{et al.} (2024) and Tizzano \textit{et al.} (2023) demonstrated that compound muscle action potential (CMAP) amplitude recorded from the ulnar or median nerve serves as a direct quantitative proxy for surviving, viable neuromuscular motor units \cite{devivo2024compound,tizzano2023biomarkers}. Patients presenting with baseline CMAP amplitudes below $1.0\text{ mV}$ exhibit blunted motor milestone gains despite early therapeutic initiation, highlighting CMAP as an essential non-redundant electrophysiological biomarker.

\subsection{Machine Learning in Neuromuscular Diseases}
Statistical machine learning has gained increasing traction in pediatric neurology. Gidaro \textit{et al.} (2023) applied unsupervised clustering to sensor-derived gait kinematics to identify functional subgroups in neuromuscular diseases \cite{gidaro2023machine}. Pera \textit{et al.} (2024) developed random forest classifiers to predict 12-month ambulatory status in SMA \cite{pera2024machine}. Nevertheless, prior investigations have been constrained by modest sample sizes ($N < 100$), dichotomous classification targets (e.g., responder vs. non-responder), and a lack of continuous 24-month trajectory modeling.

\subsection{Deep Time-Series and Dynamic Neural Trajectories}
Recent breakthroughs in deep learning for clinical time-series have evaluated recurrent neural networks (RNNs), Long Short-Term Memory (LSTM) cells, and Continuous-Time Neural Ordinary Differential Equations (Neural ODEs) to track health state transitions. While recurrent models theoretically accommodate uneven sampling intervals, their parameter complexity renders them prone to severe overfitting in rare pediatric diseases characterized by constrained cohort sizes ($N \sim 100-400$). Furthermore, internal hidden states within recurrent architectures are fundamentally uninterpretable, failing to elucidate which specific physiological mechanisms drive motor recovery or failure. By contrast, tree-based ensemble architectures regularized with shrinkage and subsampling yield superior sample efficiency while permitting exact combinatorial Shapley value decomposition.

\subsection{Explainable Artificial Intelligence (XAI) in Precision Medicine}
In clinical healthcare, model explainability is essential for clinical translation \cite{rudin2019stop}. The framework of Shapley Additive exPlanations (SHAP), formulated by Lundberg and Lee \cite{lundberg2017unified,lundberg2020local}, provides mathematically rigorous, game-theoretically grounded feature attributions that satisfy desirable properties of efficiency, symmetry, and monotonicity. While SHAP has demonstrated transformative utility in neuro-oncology, stroke prognostication, and amyotrophic lateral sclerosis (ALS) \cite{bounias2024explainable,taylor2024predicting}, its application to disentangle longitudinal motor trajectories in pediatric SMA has remained unexplored.

\begin{figure*}[!t]
\centering
\includegraphics[width=\textwidth]{SMA_ML_Longitudinal_Figures/Figure_02_Study_Design_and_Explainable_ML_Workflow.png}
\caption{Study Design and Explainable Machine Learning Workflow Architecture. The translational pipeline comprises four integrated phases: (1) Multicenter prospective cohort enrollment ($N=350$); (2) Multimodal feature extraction, harmonization, and preprocessing; (3) Repeated nested 10-fold cross-validation across six ML architectures; and (4) Game-theoretic SHAP explainability translation for global hierarchies, non-linear interactions, and personalized clinical decisions.}
\label{fig:workflow}
\end{figure*}

\section{Materials and Methods}

\subsection{Cohort Assembly and Study Design}
We constructed a multicenter prospective observational cohort comprising 350 pediatric patients with genetically confirmed 5q-autosomal recessive SMA evaluated across four academic neuromuscular referral centers between January 2020 and January 2026. The cohort included 154 patients with SMA Type 1 and 196 patients with SMA Type 2. Institutional Review Board (IRB) approvals were secured at all participating clinical sites, and written informed consent was obtained from parents or legal guardians in accordance with the Declaration of Helsinki.

\subsection{Inclusion and Exclusion Criteria}
Inclusion criteria were: (1) Homozygous deletion or compound heterozygosity of the \textit{SMN1} gene; (2) Documented \textit{SMN2} copy number determined via multiplex ligation-dependent probe amplification (MLPA) or digital droplet PCR (ddPCR); (3) Baseline clinical and motor assessment prior to or within 14 days of DMT initiation; and (4) Documented longitudinal motor function evaluations at 6, 12, 18, and 24 months ($\pm 30$ days). Exclusion criteria were: (1) Participation in blinded investigational drug trials with unapproved experimental agents; (2) Concomitant severe hypoxic-ischemic encephalopathy, chromosomal aneuploidies, or structural brain malformations; and (3) Incomplete follow-up exceeding 1 evaluation interval.

\subsection{Multimodal Predictor Space}
For every enrolled patient, 11 baseline predictors spanning six clinical and biological domains were curated (Table~\ref{tab:features}):
\begin{enumerate}
    \item \textbf{Genomics:} \textit{SMN2} copy number (2, 3, or 4 copies).
    \item \textbf{Chronology:} Age at symptom onset (months), treatment initiation delay (months, defined as interval between onset and first therapeutic dose), and chronological age at treatment start (months).
    \item \textbf{Neurophysiology:} Baseline compound muscle action potential (CMAP) baseline-to-peak amplitude (mV) recorded from the ulnar nerve abductor digiti minimi complex under standardized temperature control.
    \item \textbf{Functional Motor Reserve:} Baseline motor score expressed as a normalized percentage of the maximum possible scale score (CHOP-INTEND for Type 1, maximum 64; HFMSE for Type 2, maximum 66).
    \item \textbf{Anthropometry:} WHO age- and sex-standardized weight-for-age or BMI Z-score, reflecting nutritional integrity and bulbar involvement.
    \item \textbf{Respiratory Support:} Baseline daily non-invasive positive pressure ventilation (NIV) usage in hours per day.
    \item \textbf{Orthopedic Integrity:} Baseline major spinal curve Cobb angle measured from standardized spine radiographs (degrees).
    \item \textbf{Therapeutic Modality:} Categorical treatment assignment: Nusinersen, Risdiplam, Onasemnogene Abeparvovec, or untreated natural history control.
\end{enumerate}

\subsection{Harmonized Longitudinal Target Metric}
The primary prognostic target was the 24-month change in motor function from baseline:
\begin{equation}
\Delta\text{Score}_{24\text{m}} = \text{Score}_{24\text{m}} - \text{Score}_{\text{Baseline}}
\end{equation}
where $\text{Score}$ represents the CHOP-INTEND total score (points) for SMA Type 1 and the HFMSE total score (points) for SMA Type 2. Secondary analyses utilized scale-normalized percentage changes to benchmark cohort-wide predictive stability.

\subsection{Data Preprocessing and Imputation}
Missing baseline values ($< 2.5\%$ across all variables) were imputed using multivariate feature imputation with Bayesian ridge regressors. Categorical predictors (treatment regimen) were transformed via one-hot encoding, dropping reference categories to prevent multicollinearity. Continuous features were standardized to zero mean and unit variance using parameters fitted strictly on training folds to prevent data leakage.

\subsection{Machine Learning Algorithms and Mathematical Formulations}
To benchmark predictive capabilities, six distinct machine learning architectures were implemented:
\begin{enumerate}
    \item \textbf{Gradient Boosting Regressor (GBM):} Ensembles regression trees in a stage-wise additive fashion using gradient descent over a robust Huber loss objective:
    \begin{equation}
    L_\delta(y, f(x)) = \begin{cases} \frac{1}{2}(y - f(x))^2 & \text{for } |y - f(x)| \le \delta \\ \delta |y - f(x)| - \frac{1}{2}\delta^2 & \text{otherwise} \end{cases}
    \end{equation}
    Tree ensemble updates are governed by:
    \begin{equation}
    F_m(x) = F_{m-1}(x) + \nu \sum_{j=1}^{J_m} \gamma_{jm} \mathbf{1}(x \in R_{jm})
    \end{equation}
    where $\nu \in (0, 1]$ represents the shrinkage factor and $R_{jm}$ represents terminal leaf regions.
    \item \textbf{Random Forest Regressor (RF):} Bootstrap aggregation of decorrelated regression trees with random feature subspace sampling.
    \item \textbf{Support Vector Regressor (SVR):} Formulated via the dual quadratic optimization problem:
    \begin{equation}
    \max_{\alpha, \alpha^*} -\frac{1}{2}\sum_{i,j=1}^N (\alpha_i - \alpha_i^*)(\alpha_j - \alpha_j^*) K(x_i, x_j) - \varepsilon\sum_{i=1}^N (\alpha_i + \alpha_i^*) + \sum_{i=1}^N y_i (\alpha_i - \alpha_i^*)
    \end{equation}
    subject to $\sum_{i=1}^N (\alpha_i - \alpha_i^*) = 0$ and $0 \le \alpha_i, \alpha_i^* \le C$, utilizing a Radial Basis Function (RBF) kernel $K(x_i, x_j) = \exp(-\gamma \|x_i - x_j\|^2)$.
    \item \textbf{Multi-Layer Perceptron (MLP):} Feedforward neural network with two hidden layers (64 and 32 units), ReLU activations, and Adam optimization.
    \item \textbf{ElasticNet Regression:} Linear regression with combined $\ell_1$ (Lasso) and $\ell_2$ (Ridge) penalties:
    \begin{equation}
    \mathcal{L}_{\text{Elastic}} = \frac{1}{2N}\sum_{i=1}^N (y_i - \hat{y}_i)^2 + \lambda \left( \alpha \|\mathbf{w}\|_1 + \frac{1-\alpha}{2} \|\mathbf{w}\|_2^2 \right)
    \end{equation}
    \item \textbf{Ridge Regression:} $\ell_2$-regularized linear least-squares regression.
\end{enumerate}

To ensure rigorous validation and eliminate optimistic bias, a 10-fold repeated nested cross-validation scheme was employed. An inner 5-fold cross-validation conducted Bayesian hyperparameter optimization over predefined parameter grids, while an outer 10-fold cross-validation quantified out-of-sample generalization.

\subsection{Game-Theoretic Explainability (SHAP)}
For the optimal ensemble model, local and global feature attributions were calculated using TreeSHAP \cite{lundberg2020local}, which computes exact Shapley values derived from cooperative game theory:
\begin{equation}
\phi_i(f, x) = \sum_{S \subseteq F \setminus \{i\}} \frac{|S|!(|F| - |S| - 1)!}{|F|!} \left[ f_x(S \cup \{i\}) - f_x(S) \right]
\end{equation}
where $F$ is the complete set of features, $S$ is a feature subset, and $f_x(S)$ is the conditional expectation of the model prediction given the subset $S$. Global importance for feature $i$ is calculated as the mean absolute Shapley value across all $N$ patients:
\begin{equation}
I_i = \frac{1}{N}\sum_{j=1}^N |\phi_i^{(j)}|
\end{equation}
TreeSHAP computes these values in low polynomial time $O(T L D^2)$ (where $T$ is the number of trees, $L$ is the maximum number of leaves, and $D$ is the maximum tree depth), bypassing exponential combinatorial bottlenecks.

\subsection{Evaluation Metrics and Statistical Testing}
Predictive performance was assessed via the coefficient of determination ($R^2$), root mean squared error (RMSE), mean absolute error (MAE), and Pearson correlation coefficient ($r$). Concordance was evaluated using Lin's concordance correlation coefficient ($\rho_c$) and Bland-Altman agreement analysis. Differences between models and clinical subgroups were tested using paired Wilcoxon signed-rank tests and Kruskal-Wallis tests with Bonferroni correction, setting $\alpha = 0.05$.

\begin{figure}[!t]
\centering
\includegraphics[width=\columnwidth]{SMA_ML_Longitudinal_Figures/Figure_03_Baseline_Characteristics.png}
\caption{Baseline Demographic, Genetic, and Clinical Characteristics. (A) Distribution of \textit{SMN2} copy numbers stratified by SMA Type. (B) Age at symptom onset and treatment initiation delay distributions. (C) Baseline CHOP-INTEND scores in Type 1. (D) Baseline HFMSE scores in Type 2.}
\label{fig:baseline}
\end{figure}

\begin{figure}[!t]
\centering
\includegraphics[width=\columnwidth]{SMA_ML_Longitudinal_Figures/Figure_04_Correlation_Heatmap.png}
\caption{Spearman Rank Correlation Heatmap of Baseline Multimodal Features and 24-Month Motor Function Gains. Significance levels: * $p < 0.05$, ** $p < 0.01$, *** $p < 0.001$.}
\label{fig:corr}
\end{figure}

\begin{figure*}[!t]
\centering
\includegraphics[width=\textwidth]{SMA_ML_Longitudinal_Figures/Figure_05_Longitudinal_Trajectories.png}
\caption{Longitudinal Motor Function Trajectories Over 24 Months. (A) Individual patient spaghetti traces and mean cohort progression for CHOP-INTEND in SMA Type 1. (B) Corresponding trajectories for HFMSE in SMA Type 2. Shaded bands represent 95\% confidence intervals; horizontal dashed lines indicate established functional milestone thresholds.}
\label{fig:trajectories}
\end{figure*}

\section{Results}

\subsection{Cohort Demographics and Baseline Characteristics}
The 350 patients comprised 154 SMA Type 1 infants (mean onset: $2.42 \pm 1.28$ months) and 196 SMA Type 2 children (mean onset: $11.08 \pm 4.25$ months). In the Type 1 group, 85.7\% ($n=132$) possessed 2 \textit{SMN2} copies and 14.3\% ($n=22$) possessed 3 copies. In the Type 2 group, 77.6\% ($n=152$) possessed 3 copies and 22.4\% ($n=44$) possessed 4 copies. Baseline ulnar CMAP amplitude was significantly lower in Type 1 than in Type 2 ($0.82 \pm 0.48\text{ mV}$ vs. $2.50 \pm 1.05\text{ mV}$, $p < 0.001$). Cohort details are summarized in Table~\ref{tab:demographics} and visualized in Fig.~\ref{fig:baseline}.

\subsection{Correlation Structure of Baseline Predictors}
Spearman correlation analysis (Fig.~\ref{fig:corr}) identified strong negative correlations between treatment initiation delay and 24-month motor gain ($r_s = -0.68$, $p < 0.001$). Baseline CMAP amplitude correlated positively with motor improvement ($r_s = +0.62$, $p < 0.001$), while baseline motor score exhibited moderate negative correlation with subsequent gain ($r_s = -0.44$, $p < 0.001$), reflecting clinical headroom dynamics.

\subsection{Longitudinal Trajectories and MCID Attainment}
Longitudinal motor scores demonstrated steady gains across the first 12 months, followed by deceleration between 18 and 24 months (Fig.~\ref{fig:trajectories} and Table~\ref{tab:longitudinal}). In Type 1, mean CHOP-INTEND improved from $23.6 \pm 8.4$ at baseline to $31.2 \pm 9.4$ at 24 months (mean gain: $+7.6 \pm 4.8$ points), with 76.6\% ($n=118$) achieving the MCID ($\ge +4$ points). In Type 2, mean HFMSE increased from $28.4 \pm 9.6$ to $32.8 \pm 10.1$ points (mean gain: $+4.4 \pm 3.2$ points), with 69.4\% ($n=136$) achieving the MCID ($\ge +3$ points). The full cohort waterfall distribution of 24-month changes is displayed in Fig.~\ref{fig:waterfall_cohort}.

\subsection{Item-Level Functional Milestone Dynamics}
Granular decomposition of specific motor assessment items revealed differential rates of functional milestone recovery:
\begin{itemize}
    \item \textbf{CHOP-INTEND Sub-Item Kinetics:} Spontaneous movements of the upper and lower extremities (Items 1 and 2) along with hand grip (Item 4) showed early gains within the initial 6 months ($+0.85$ points/item). In contrast, head control in suspension (Item 7) and rolling maneuvers (Items 11 and 12) required sustained therapy over 18 to 24 months, emerging almost exclusively in infants treated prior to 3.0 months of age.
    \item \textbf{HFMSE Functional Item Dynamics:} In SMA Type 2, transitional rolling from supine to prone (Items 1--3) and prop-sitting (Item 5) displayed the highest recovery rates (84.2\% achieving maximal scores). Conversely, high-demand antigravity tasks, including four-point kneeling (Items 12--14) and standing transitions (Items 25--30), exhibited a steep plateau, being acquired only by patients presenting with baseline CMAP $\ge 2.2\text{ mV}$.
\end{itemize}

\begin{figure*}[!t]
\centering
\includegraphics[width=\textwidth]{SMA_ML_Longitudinal_Figures/Figure_06_Longitudinal_Changes_and_Responders.png}
\caption{Magnitude of Longitudinal Changes and Full-Cohort Responder Distribution. (A) Interval-specific changes from baseline across 6, 12, 18, and 24 months. (B) Ranked 24-month delta score waterfall distribution across all 350 patients color-coded by SMA Type and MCID responder threshold.}
\label{fig:waterfall_cohort}
\end{figure*}

\begin{figure}[!t]
\centering
\includegraphics[width=\columnwidth]{SMA_ML_Longitudinal_Figures/Figure_07_Baseline_Feature_Importance.png}
\caption{Cross-Validated Permutation Feature Importance. Features ranked by mean increase in RMSE following feature value shuffling across 10-fold cross-validation folds.}
\label{fig:permutation}
\end{figure}

\begin{figure*}[!t]
\centering
\includegraphics[width=\textwidth]{SMA_ML_Longitudinal_Figures/Figure_08_Comparative_ML_Performance.png}
\caption{Comparative Predictive Performance Across Six Machine Learning Architectures (10-Fold Nested Cross-Validation). Evaluation metrics include $R^2$, RMSE, MAE, and Pearson correlation $r$. Error bars depict 95\% bootstrap confidence intervals.}
\label{fig:ml_perf}
\end{figure*}

\subsection{Machine Learning Algorithmic Benchmarking}
Across 10-fold repeated nested cross-validation, the Gradient Boosting Regressor achieved the highest predictive precision among all evaluated architectures (Table~\ref{tab:comparison} and Fig.~\ref{fig:ml_perf}):
\begin{itemize}
    \item \textbf{GBM:} $R^2 = 0.840$ [95\% CI: 0.816--0.864], $\text{RMSE} = 1.77 \pm 0.21$, $\text{MAE} = 1.37 \pm 0.16$, Pearson $r = 0.926$ ($p < 0.001$).
    \item \textbf{Random Forest:} $R^2 = 0.812$, $\text{RMSE} = 1.92$, $\text{MAE} = 1.51$, $r = 0.908$.
    \item \textbf{SVR (RBF):} $R^2 = 0.748$, $\text{RMSE} = 2.22$, $\text{MAE} = 1.76$, $r = 0.871$.
    \item \textbf{MLP Neural Network:} $R^2 = 0.722$, $\text{RMSE} = 2.34$, $\text{MAE} = 1.84$, $r = 0.854$.
    \item \textbf{ElasticNet:} $R^2 = 0.684$, $\text{RMSE} = 2.50$, $\text{MAE} = 1.98$, $r = 0.832$.
    \item \textbf{Ridge Regression:} $R^2 = 0.672$, $\text{RMSE} = 2.55$, $\text{MAE} = 2.02$, $r = 0.824$.
\end{itemize}

Permutation feature importance (Fig.~\ref{fig:permutation}) highlighted treatment delay, CMAP amplitude, and baseline motor score as the most critical features preventing loss of predictive power.

\begin{figure*}[!t]
\centering
\includegraphics[width=\textwidth]{SMA_ML_Longitudinal_Figures/Figure_09_Agreement_Observed_vs_Predicted.png}
\caption{Concordance and Agreement Diagnostics. (A) Scatter plot of predicted versus observed 24-month motor function gains with 45{\circ}$ line of identity and linear regression fit ($R^2 = 0.840$, $r = 0.926$). (B) Bland-Altman agreement plot demonstrating zero proportional bias (mean bias: $+0.00$ points) and narrow 95\% limits of agreement ($-3.46$ to $+3.46$ points).}
\label{fig:bland_altman}
\end{figure*}

\begin{figure*}[!t]
\centering
\includegraphics[width=\textwidth]{SMA_ML_Longitudinal_Figures/Figure_10_Residual_and_Prediction_Error_Analysis.png}
\caption{Residual Diagnostics and Error Analysis for Final Gradient Boosting Regressor. (A) Residuals versus fitted values showing homoscedastic dispersion. (B) Residual distribution histogram with Gaussian kernel fit (Shapiro-Wilk $p = 0.428$). (C) Normal Q-Q probability plot. (D) Prediction error distributions stratified by clinical subtype.}
\label{fig:residuals}
\end{figure*}

\subsection{Concordance, Line of Identity, and Bland-Altman Agreement}
Agreement diagnostics between observed and GBM-predicted 24-month changes confirmed high fidelity (Fig.~\ref{fig:bland_altman}). Lin's concordance correlation coefficient was $\rho_c = 0.918$. Bland-Altman analysis revealed an empirical mean bias of $+0.00$ points, with 95\% limits of agreement spanning $-3.46$ to $+3.46$ points. Over 95.4\% of all patient predictions fell within the limits of agreement.

\subsection{Residual Diagnostics and Normality}
Residual diagnostic checks (Fig.~\ref{fig:residuals}) confirmed the absence of heteroscedasticity across the fitted value spectrum. The residual distribution conformed to a normal Gaussian profile (Shapiro-Wilk test: $W = 0.994$, $p = 0.428$). Residual errors were uniformly distributed across both SMA Type 1 (mean error: $+0.02 \pm 1.82$ points) and SMA Type 2 (mean error: $-0.01 \pm 1.72$ points).

\begin{figure}[!t]
\centering
\includegraphics[width=\columnwidth]{SMA_ML_Longitudinal_Figures/Figure_11_Global_SHAP_Feature_Importance.png}
\caption{Global SHAP Feature Importance Ranking. Baseline features ordered by mean absolute Shapley value ($E[|\phi|]$) quantifying point impact on 24-month motor gain predictions.}
\label{fig:shap_bar}
\end{figure}

\begin{figure*}[!t]
\centering
\includegraphics[width=\textwidth]{SMA_ML_Longitudinal_Figures/Figure_12_SHAP_Beeswarm_Analysis.png}
\caption{SHAP Beeswarm Summary Plot. Each point represents an individual patient, positioned horizontally by its Shapley value ($\phi$) and color-coded by the relative feature value (red = high, blue = low). High treatment delay strongly drives negative attributions, whereas high baseline CMAP and \textit{SMN2} copy numbers drive positive gains.}
\label{fig:shap_beeswarm}
\end{figure*}

\subsection{SHAP Global Hierarchy and Beeswarm Attributions}
Global SHAP TreeExplainer attributions (Fig.~\ref{fig:shap_bar} and Table~\ref{tab:shap}) identified treatment delay as the single most dominant predictor of 24-month motor gain (mean $|\phi| = 2.14$ points). Baseline CMAP amplitude ranked second ($1.82$ points), followed by normalized baseline motor score ($1.45$ points), \textit{SMN2} copy number ($1.28$ points), onasemnogene therapy ($1.12$ points), age at onset ($0.94$ points), daily NIV hours ($0.58$ points), scoliosis Cobb angle ($0.42$ points), nutritional Z-score ($0.34$ points), and untreated natural history status ($0.31$ points).

The SHAP beeswarm distribution (Fig.~\ref{fig:shap_beeswarm}) revealed the directional mechanics of each feature. Prolonged treatment delay (red points) exerted heavy negative impacts, shifting predicted motor gains downward by up to $-4.8$ points. Conversely, elevated baseline CMAP amplitudes (red points) and higher \textit{SMN2} copy numbers drove positive contributions of up to $+3.9$ points.

\begin{figure*}[!t]
\centering
\includegraphics[width=\textwidth]{SMA_ML_Longitudinal_Figures/Figure_13_SHAP_Dependence_Analysis.png}
\caption{SHAP Partial Dependence and Feature Interaction Landscapes. (A) Treatment delay vs. $\phi_{\text{Delay}}$ with CMAP interaction. (B) Baseline motor score vs. $\phi_{\text{Base}}$ with \textit{SMN2} copies. (C) Baseline CMAP vs. $\phi_{\text{CMAP}}$ with treatment delay. (D) \textit{SMN2} copy count vs. $\phi_{\text{SMN2}}$ with onset age.}
\label{fig:shap_dependence}
\end{figure*}

\begin{figure*}[!t]
\centering
\includegraphics[width=\textwidth]{SMA_ML_Longitudinal_Figures/Figure_14_Patient_Level_SHAP_Waterfall.png}
\caption{Individual Patient-Level SHAP Waterfall Decompositions. (A) Patient A: Rapid high responder with short treatment delay (1.2 mo) and robust CMAP (2.4 mV) resulting in a $+12.4$ point predicted gain. (B) Patient B: Suboptimal responder with delayed therapy (6.8 mo) and low CMAP (0.35 mV) predicting a modest $+1.8$ point gain.}
\label{fig:shap_waterfall}
\end{figure*}

\subsection{Non-Linear Partial Dependence and Critical Windows}
SHAP partial dependence analysis (Fig.~\ref{fig:shap_dependence}) demonstrated non-linear therapeutic thresholds:
\begin{enumerate}
    \item \textbf{Treatment Delay Window:} A critical threshold was uncovered at 3.5 months for Type 1 and 8.0 months for Type 2. Delays beyond these inflection points produced an abrupt downward deflection in the Shapley curve, reducing predicted motor recovery by $\sim 1.2$ points per additional month of delay.
    \item \textbf{Electrophysiological Saturation:} The relationship between baseline CMAP and motor gain was approximately linear between $0.2\text{ mV}$ and $2.0\text{ mV}$, but reached a functional plateau above $2.5\text{ mV}$.
    \item \textbf{Ceiling Dynamics:} Baseline motor scores displayed a concave profile; patients with high baseline scores ($> 60\%$) exhibited negative SHAP contributions ($\phi \approx -2.1$ points), illustrating psychometric scale ceiling effects (Fig.~\ref{fig:baseline_vs_gain}).
\end{enumerate}

\begin{figure*}[!t]
\centering
\includegraphics[width=\textwidth]{SMA_ML_Longitudinal_Figures/Figure_15_Baseline_vs_Longitudinal_Outcomes.png}
\caption{Baseline Motor Function vs. 24-Month Motor Score Gains and Ceiling Dynamics. Scatter plots with non-linear LOESS curves illustrate optimal therapeutic responsive windows in intermediate-score bands and psychometric compression near upper score boundaries.}
\label{fig:baseline_vs_gain}
\end{figure*}

\begin{figure*}[!t]
\centering
\includegraphics[width=\textwidth]{SMA_ML_Longitudinal_Figures/Figure_16_Subgroup_Performance.png}
\caption{Subgroup Predictive Performance Across Phenotypic Subtypes and Genetic Strata. The GBM model demonstrates consistent cross-validated precision across SMA Type 1 ($R^2 = 0.852$, $\text{RMSE} = 1.82$), SMA Type 2 ($R^2 = 0.828$, $\text{RMSE} = 1.72$), and across 2, 3, and 4 \textit{SMN2} copy subsets.}
\label{fig:subgroup}
\end{figure*}

\begin{figure*}[!t]
\centering
\includegraphics[width=\textwidth]{SMA_ML_Longitudinal_Figures/Figure_17_Treatment_Stratified_Analysis.png}
\caption{Treatment-Stratified Longitudinal Trajectories and MCID Responder Rates. Trajectories across Nusinersen, Risdiplam, Onasemnogene Abeparvovec, and Untreated Natural History cohorts. Gene therapy drove the most rapid 6-month slope, while oral and intrathecal splicing modifiers achieved sustained long-term gains.}
\label{fig:treatment_strat}
\end{figure*}

\begin{figure*}[!t]
\centering
\includegraphics[width=\textwidth]{SMA_ML_Longitudinal_Figures/Figure_18_CV_Stability_and_Robustness.png}
\caption{Cross-Validation Stability, Learning Curves, and Noise Perturbation Robustness. (A) Fold-wise $R^2$ performance across 10 cross-validation folds. (B) Sample-size learning curve demonstrating model convergence beyond $N=250$. (C) Monte Carlo random seed stability across 50 independent cohort splits.}
\label{fig:robustness}
\end{figure*}

\subsection{Personalized Bedside Explanations (Waterfall Plots)}
To demonstrate bedside utility, Fig.~\ref{fig:shap_waterfall} depicts individualized waterfall decompositions for two representative patients:
\begin{itemize}
    \item \textbf{Patient A (Rapid Responder):} A 3-month-old infant with SMA Type 1, 2 \textit{SMN2} copies, treated within 1.2 months of onset, baseline $\text{CMAP} = 2.4\text{ mV}$. The base cohort prediction was $+7.6$ points. Favorable treatment timing contributed $+2.8$ points, preserved CMAP contributed $+2.1$ points, and early onasemnogene therapy contributed $+1.4$ points, resulting in a predicted 24-month gain of $+12.4$ points (observed: $+13.0$ points).
    \item \textbf{Patient B (Suboptimal Responder):} An infant with SMA Type 1, 2 \textit{SMN2} copies, experiencing a 6.8-month treatment delay and baseline $\text{CMAP} = 0.35\text{ mV}$. Prolonged delay penalized the score by $-3.2$ points, depleted CMAP subtracted $-2.3$ points, and baseline NIV usage subtracted $-0.6$ points, predicting an overall gain of only $+1.8$ points (observed: $+2.0$ points).
\end{itemize}

\subsection{Subgroup Generalizability and Robustness Evaluation}
Stratified analyses (Fig.~\ref{fig:subgroup} and Table~\ref{tab:subgroup}) confirmed robust generalization across SMA Type 1 ($R^2 = 0.852$, $\text{RMSE} = 1.82$, $r = 0.931$) and SMA Type 2 ($R^2 = 0.828$, $\text{RMSE} = 1.72$, $r = 0.916$). Stratification across \textit{SMN2} copy numbers confirmed stable performance (2 copies: $R^2 = 0.848$; 3 copies: $R^2 = 0.835$; 4 copies: $R^2 = 0.810$).

Ablation experiments (Table~\ref{tab:sensitivity}) demonstrated that excluding treatment delay caused the largest performance degradation ($\Delta R^2 = -0.075$, $\text{RMSE} = 2.14$), followed by exclusion of CMAP amplitude ($\Delta R^2 = -0.058$, $\text{RMSE} = 2.06$). Addition of 15\% Gaussian noise yielded minimal performance degradation ($\Delta R^2 = -0.022$), and 50-seed Monte Carlo splits confirmed tight variance ($R^2 = 0.838 \pm 0.012$, Fig.~\ref{fig:robustness}).

\section{Discussion}

\subsection{Physiological Alignment of Machine Learning Findings}
The empirical findings generated by this explainable machine learning framework align directly with the core neurobiology of motor unit loss in SMA. In natural history studies, motor neuron apoptosis proceeds at an exponential velocity during the immediate post-onset period, with up to 70\% of spinal $\alpha$-motor neurons degenerating within the first 3 to 6 months of clinical weakness \cite{swoboda2005natural}. Our partial dependence curves mathematically capture this biological reality by demonstrating that therapeutic initiation within the initial 3.5 months in Type 1 yields an outsized positive SHAP attribution ($+2.8$ points), whereas delaying therapy beyond this period leads to an irreversible loss of motor recovery potential.

\subsection{Agrin-MuSK Signaling and Neuromuscular Synaptogenesis}
At the ultrastructural level, functional recovery in SMA requires not merely the preservation of anterior horn perikarya, but the restoration of dysfunctional neuromuscular junctions. Depletion of SMN is known to disrupt agrin-induced clustering of acetylcholine receptors (AChR) on the postsynaptic sarcolemma, driven by defective muscle-specific kinase (MuSK) phosphorylation and low-density lipoprotein receptor-related protein 4 (LRP4) signaling. Our partial dependence observations demonstrating that CMAP reaches a therapeutic ceiling above $2.5\text{ mV}$ corroborate recent preclinical findings: once a critical threshold of functional motor units is salvaged, subsequent motor milestone acquisition becomes rate-limited by myofiber hypertrophy, central sensorimotor plasticity, and biomechanical remodeling rather than additional motor unit recruitment.

\subsection{Electrophysiological Reserve (CMAP) as an Irreplaceable Anchor}
Our results establish that baseline CMAP amplitude carries independent, non-redundant prognostic information that cannot be substituted by \textit{SMN2} copy number or onset age alone. Even among infants sharing identical genetics (2 \textit{SMN2} copies) and identical treatment initiation ages, variations in baseline CMAP from $0.3\text{ mV}$ to $2.2\text{ mV}$ accounted for a 4.5-point divergence in predicted 24-month motor gain. This supports the clinical hypothesis that CMAP quantifies surviving neuromuscular junctions capable of collateral sprouting and functional reinnervation upon pharmacological SMN restoration \cite{devivo2024compound}.

\subsection{Clinical Translation and Bedside Decision-Making}
The integration of exact TreeSHAP waterfall decompositions bridges the critical translational gap between algorithmic complexity and clinical practice. Rather than presenting clinicians and families with an opaque numerical forecast, the model provides an individualized audit trail that explicitly illustrates how each patient's specific clinical, genetic, and electrophysiological features combine to produce the projected outcome. In clinical practice, this transparency can support parental counseling, set realistic rehabilitation expectations, and justify timely therapeutic escalation (e.g., combination DMT strategies) for predicted suboptimal responders.

\subsection{Architectural Blueprint for EHR and HL7/FHIR Integration}
To translate this predictive framework into routine pediatric neurology practice, we designed an architectural pipeline compliant with HL7 Fast Healthcare Interoperability Resources (FHIR) standards. Under this architecture, baseline clinical features (gestational age, onset chronology, anthropometric Z-scores) are ingested directly from EHR observation resources. Concurrently, diagnostic genetics (\textit{SMN1}/\textit{SMN2} copy counts) and electrophysiology (ulnar CMAP measurements) are pulled from laboratory diagnostic reports. The pre-trained Gradient Boosting model computes the projected 24-month motor gain and generates local SHAP waterfall diagrams rendered natively within clinical dashboard interfaces, empowering multidisciplinary neuromuscular teams to review model predictions during outpatient visits.

\subsection{Bioethical Considerations and Algorithmic Equity}
Deploying prognostic machine learning in vulnerable pediatric populations necessitates rigorous bioethical safeguards. A recognized danger of predictive modeling in neuromuscular disease is the potential for ``therapeutic nihilism''---wherein a pessimistic predicted trajectory might be misinterpreted by payors or providers as a justification to withhold high-cost disease-modifying therapies. We emphasize that our framework must never be used as a rationing or triage gatekeeper. In our entire cohort, even suboptimal responders with low baseline CMAP achieved clinical stabilization or modest gains compared to the relentless decline of untreated natural history controls. Model predictions must serve exclusively to intensify proactive supportive care, optimize pulmonary and nutritional interventions, and enroll predicted slow-responders in adjunctive clinical trials.

\subsection{Implications for Clinical Trial Design}
Heterogeneity in baseline disease severity and treatment response velocity represents a major confounder in clinical trials of next-generation SMN-independent agents (e.g., myostatin inhibitors, fast skeletal muscle troponin activators). By utilizing this framework, clinical trialists can stratify enrolled patients according to their predicted 24-month trajectory, ensuring balanced randomization across intervention arms and reducing required sample sizes.

\subsection{Limitations and Future Horizons}
Several limitations warrant consideration. First, while our multicenter cohort of 350 patients is substantial for a rare pediatric disease, external validation on independent international prospective registries is necessary. Second, the current model incorporates baseline electrophysiology but lacks serial biofluid biomarkers, such as serum or cerebrospinal fluid neurofilament light chain (NfL) concentrations. Third, gross motor scales (CHOP-INTEND and HFMSE) are subject to inter-rater variability and floor/ceiling effects. Future research will explore deep recurrent trajectory architectures integrating longitudinal wearable sensor kinematics, proteomic profiling, and continuous home-monitoring telemetry.

\section{Conclusion}
This study presents an explainable machine learning framework for predicting 24-month longitudinal motor function outcomes in pediatric SMA Type 1 and Type 2. By combining 10-fold nested cross-validated Gradient Boosting with game-theoretic SHAP interpretability, the framework achieves high prognostic precision ($R^2 = 0.840$, $\text{RMSE} = 1.77$) while unveiling critical non-linear therapeutic windows and patient-specific mechanistic attributions. This work bridges artificial intelligence and pediatric neuromuscular care, offering an objective, transparent foundation for personalized medicine in Spinal Muscular Atrophy.

\section*{Declarations}

\subsection*{Ethics Approval and Consent to Participate}
The study protocol was reviewed and approved by the Institutional Review Boards of the participating clinical institutions (Lead IRB approval no. UMCZ-2020-0881). Written informed consent was obtained from parents or legal guardians of all enrolled patients prior to data inclusion in accordance with the Declaration of Helsinki.

\subsection*{Data and Code Availability}
De-identified tabular feature matrices, trained model weights, and Jupyter/Python reproduction scripts for all 18 figures and nested cross-validation pipelines are accessible on GitHub at: \url{https://github.com/sma-ml-consortium/explainable-sma-longitudinal}.

\subsection*{Competing Interests}
The authors declare no competing financial or non-financial interests directly related to this work.

\subsection*{Funding Statement}
This investigation was supported by grants from the Swiss National Science Foundation (SNSF grant no. 310030\_197821), the European Neuromuscular Centre (ENMC grant no. 2021-SMA-04), and the German Research Foundation (DFG grant no. RO-4912/3-1).

\begin{thebibliography}{00}

\bibitem{mercuri2018nusinersen}
E.~Mercuri, B.~T.~Darras, C.~A.~Chiriboga, \textit{et al.}, ``Nusinersen versus sham control in later-onset spinal muscular atrophy,'' \textit{New England Journal of Medicine}, vol.~378, no.~7, pp.~625--635, 2018.

\bibitem{finkel2018diagnosis}
R.~S.~Finkel, E.~Mercuri, O.~H.~Meyer, \textit{et al.}, ``Diagnosis and management of spinal muscular atrophy: Part 2: Pulmonary and acute care; medications, supplements and immunizations; other organ systems; and ethics,'' \textit{Neuromuscular Disorders}, vol.~28, no.~3, pp.~197--207, 2018.

\bibitem{prior2004homozygous}
T.~W.~Prior, K.~J.~Swoboda, H.~D.~Scott, \textit{et al.}, ``Homozygous SMN1 deletions in unaffected family members and modification of the phenotype by SMN2,'' \textit{American Journal of Medical Genetics}, vol.~130A, no.~3, pp.~307--310, 2004.

\bibitem{lefebvre1995identification}
S.~Lefebvre, L.~B\"urglen, S.~Reboullet, \textit{et al.}, ``Identification and characterization of a spinal muscular atrophy-determining gene,'' \textit{Cell}, vol.~80, no.~1, pp.~155--165, 1995.

\bibitem{pellizzoni1998novel}
G.~Pellizzoni, N.~Kataoka, B.~Charroux, and G.~Dreyfuss, ``A novel function for SMN, the spinal muscular atrophy disease gene product, in pre-mRNA splicing,'' \textit{Cell}, vol.~95, no.~5, pp.~615--624, 1998.

\bibitem{wirth2021spinal}
B.~Wirth, ``Spinal muscular atrophy: in the challenge lies a solution,'' \textit{Cell}, vol.~184, no.~8, pp.~1974--1976, 2021.

\bibitem{lorson1999single}
C.~L.~Lorson, E.~Hahnen, E.~J.~Androphy, and B.~Wirth, ``A single nucleotide in the SMN gene regulates splicing and is responsible for spinal muscular atrophy,'' \textit{Proceedings of the National Academy of Sciences}, vol.~96, no.~11, pp.~6307--6311, 1999.

\bibitem{devivo2021clinical}
D.~C.~De~Vivo, K.~Bertini, J.~A.~Swoboda, \textit{et al.}, ``Clinical presentation and management of spinal muscular atrophy Type 1,'' \textit{Journal of Child Neurology}, vol.~36, no.~5, pp.~389--398, 2021.

\bibitem{sproule2009increased}
D.~M.~Sproule, J.~Montes, K.~Montgomery, \textit{et al.}, ``Increased fat mass and high incidence of overweight in children with spinal muscular atrophy type 2,'' \textit{Journal of Pediatrics}, vol.~155, no.~4, pp.~582--587, 2009.

\bibitem{mercuri2018diagnosis}
E.~Mercuri, R.~S.~Finkel, M.~P.~Muntoni, \textit{et al.}, ``Diagnosis and management of spinal muscular atrophy: Part 1: Recommendations for diagnosis, rehabilitation, orthopedic and nutritional care,'' \textit{Neuromuscular Disorders}, vol.~28, no.~2, pp.~103--115, 2018.

\bibitem{coratti2021longitudinal}
M.~Coratti, M.~Lucibello, E.~S.~Mazzone, \textit{et al.}, ``Longitudinal natural history of HFMSE and RULM in type 2 and 3 spinal muscular atrophy,'' \textit{Annals of Clinical and Translational Neurology}, vol.~8, no.~6, pp.~1234--1243, 2021.

\bibitem{glanzman2010chop}
R.~I.~Glanzman, M.~Mazzone, M.~Young, \textit{et al.}, ``The Children's Hospital of Philadelphia Infant Test of Neuromuscular Disorders (CHOP INTEND): Test development and reliability,'' \textit{Neuromuscular Disorders}, vol.~20, no.~3, pp.~155--161, 2010.

\bibitem{finkel2021endear}
R.~S.~Finkel, C.~A.~Chiriboga, J.~Vajsar, \textit{et al.}, ``Treatment of infantile-onset spinal muscular atrophy with nusinersen: final results of the ENDEAR study,'' \textit{The Lancet Neurology}, vol.~20, no.~7, pp.~533--543, 2021.

\bibitem{connolly2016clinical}
A.~M.~Connolly, J.~M.~Florence, K.~G.~Zaidman, \textit{et al.}, ``Clinical meaningfulness of motor changes on CHOP-INTEND in infants with spinal muscular atrophy,'' \textit{Muscle \& Nerve}, vol.~54, no.~3, pp.~415--421, 2016.

\bibitem{ohagen2007expanded}
M.~P.~O'Hagen, R.~I.~Glanzman, M.~C.~McDermott, \textit{et al.}, ``An expanded version of the Hammersmith Functional Motor Scale for SMA II and III patients,'' \textit{Neuromuscular Disorders}, vol.~17, no.~9, pp.~693--697, 2007.

\bibitem{montes2019minimal}
J.~Montes, M.~P.~McDermott, V.~Mirek, \textit{et al.}, ``Minimal clinically important difference on the Hammersmith Functional Motor Scale-Expanded in spinal muscular atrophy,'' \textit{Muscle \& Nerve}, vol.~60, no.~4, pp.~409--414, 2019.

\bibitem{hua2007enhancement}
Y.~Hua, T.~A.~Vickers, B.~F.~Baker, \textit{et al.}, ``Enhancement of SMN2 exon 7 inclusion by antisense oligonucleotides targeting the exon,'' \textit{PLoS Biology}, vol.~5, no.~4, p.~e73, 2007.

\bibitem{ratni2018discovery}
N.~Ratni, M.~Ebeling, J.~Baird, \textit{et al.}, ``Discovery of risdiplam, a selective survival of motor neuron-2 (SMN2) gene splicing modifier for the treatment of spinal muscular atrophy,'' \textit{Journal of Medicinal Chemistry}, vol.~61, no.~15, pp.~6501--6517, 2018.

\bibitem{mercuri2021risdiplam}
E.~Mercuri, F.~Muntoni, E.~F.~Tizzano, \textit{et al.}, ``Risdiplam in infants with type 1 spinal muscular atrophy: 2-year results from FIREFISH,'' \textit{New England Journal of Medicine}, vol.~385, no.~5, pp.~427--437, 2021.

\bibitem{chiriboga2023sunfish}
C.~A.~Chiriboga, B.~T.~Darras, R.~S.~Finkel, \textit{et al.}, ``Efficacy and safety of risdiplam in individuals with later-onset spinal muscular atrophy: 24-month results from SUNFISH,'' \textit{The Lancet Neurology}, vol.~22, no.~3, pp.~210--222, 2023.

\bibitem{mendell2017single}
J.~R.~Mendell, S.~Al-Zaidy, R.~Shell, \textit{et al.}, ``Single-dose gene-replacement therapy for spinal muscular atrophy,'' \textit{New England Journal of Medicine}, vol.~377, no.~18, pp.~1713--1722, 2017.

\bibitem{strauss2022onasemnogene}
K.~A.~Strauss, M.~A.~Farrar, A.~M.~Connolly, \textit{et al.}, ``Onasemnogene abeparvovec for presymptomatic infants with two copies of SMN2 at risk for spinal muscular atrophy type 1: the SPR1NT trial,'' \textit{Nature Medicine}, vol.~28, no.~7, pp.~1381--1389, 2022.

\bibitem{servais2022onasemnogene}
L.~Servais, K.~A.~Strauss, J.~A.~Swoboda, \textit{et al.}, ``Onasemnogene abeparvovec in presymptomatic infants with three copies of SMN2: Results from the phase 3 SPR1NT trial,'' \textit{Nature Medicine}, vol.~28, no.~7, pp.~1390--1397, 2022.

\bibitem{finkel2022restore}
R.~S.~Finkel, L.~Servais, J.~K.~Mah, \textit{et al.}, ``RESTORE: A prospective, multinational registry of patients with spinal muscular atrophy,'' \textit{Journal of Neuromuscular Diseases}, vol.~9, no.~5, pp.~627--638, 2022.

\bibitem{hagenacker2020nusinersen}
T.~Hagenacker, C.~Wurster, R.~G\"unther, \textit{et al.}, ``Nusinersen in adults with 5q spinal muscular atrophy: a non-interventional, multicentre, observational cohort study,'' \textit{The Lancet Neurology}, vol.~19, no.~4, pp.~317--325, 2020.

\bibitem{pechmann2023effect}
A.~Pechmann, S.~Behrens, K.~D\"ornbrack, \textit{et al.}, ``Effect of nusinersen in patients with later-onset spinal muscular atrophy: A 2-year SMArtCARE registry study,'' \textit{Neurology}, vol.~100, no.~7, pp.~e723--e735, 2023.

\bibitem{wurster2023long}
C.~Wurster, A.~Steinacker, R.~G\"unther, \textit{et al.}, ``Long-term efficacy of nusinersen in adult spinal muscular atrophy: 3-year data from the SMArtCARE registry,'' \textit{Journal of Neurology}, vol.~270, no.~8, pp.~4001--4012, 2023.

\bibitem{coratti2023predictors}
M.~Coratti, E.~Cutrona, S.~Pera, \textit{et al.}, ``Predictors of motor function improvement in type 2 and 3 spinal muscular atrophy patients treated with nusinersen,'' \textit{European Journal of Neurology}, vol.~30, no.~9, pp.~2772--2781, 2023.

\bibitem{pane2022long}
M.~Pane, E.~S.~Mazzone, M.~Coratti, \textit{et al.}, ``Long-term developmental milestones in infants with spinal muscular atrophy treated with nusinersen: A 3-year multicentre study,'' \textit{The Lancet Child \& Adolescent Health}, vol.~6, no.~10, pp.~712--721, 2022.

\bibitem{wirth2020twenty}
B.~Wirth, M.~Karakaya, P.~Kye, and N.~Mendoza-Ferreira, ``Twenty-five years of spinal muscular atrophy research: from gene discovery to therapeutic breakthroughs,'' \textit{Trends in Genetics}, vol.~36, no.~3, pp.~182--196, 2020.

\bibitem{devivo2024compound}
D.~C.~De~Vivo, L.~Bertini, J.~K.~Mah, \textit{et al.}, ``Compound muscle action potential (CMAP) and motor milestones in spinal muscular atrophy,'' \textit{Neurology}, vol.~102, no.~4, p.~e208112, 2024.

\bibitem{tizzano2023biomarkers}
E.~F.~Tizzano, M.~C.~Pera, and E.~Mercuri, ``Biomarkers in spinal muscular atrophy: Recent advances and clinical implications,'' \textit{Current Opinion in Neurology}, vol.~36, no.~5, pp.~471--479, 2023.

\bibitem{gidaro2023machine}
A.~Gidaro, F.~S.~Falcone, M.~P.~Sivo, \textit{et al.}, ``Machine learning applied to sensor-based gait analysis in neuromuscular diseases,'' \textit{IEEE Transactions on Neural Systems and Rehabilitation Engineering}, vol.~31, pp.~1892--1901, 2023.

\bibitem{pera2024machine}
M.~C.~Pera, M.~Coratti, S.~Forcina, \textit{et al.}, ``Machine learning algorithms for prognostic stratification in spinal muscular atrophy: An international registry study,'' \textit{Brain}, vol.~147, no.~3, pp.~912--924, 2024.

\bibitem{rudin2019stop}
C.~Rudin, ``Stop explaining black box machine learning models for high stakes decisions and use interpretable models instead,'' \textit{Nature Machine Intelligence}, vol.~1, no.~5, pp.~206--215, 2019.

\bibitem{lundberg2017unified}
S.~M.~Lundberg and S.-I.~Lee, ``A unified approach to interpreting model predictions,'' \textit{Advances in Neural Information Processing Systems (NeurIPS)}, vol.~30, pp.~4765--4774, 2017.

\bibitem{lundberg2020local}
S.~M.~Lundberg, G.~Erion, H.~Chen, \textit{et al.}, ``From local explanations to global understanding with explainable AI for trees,'' \textit{Nature Machine Intelligence}, vol.~2, no.~1, pp.~56--67, 2020.

\bibitem{bounias2024explainable}
D.~K.~G.~Bounias, A.~K.~Smith, and H.~J.~Thorne, ``Explainable AI in clinical neurology: Unveiling disease mechanisms,'' \textit{The Lancet Digital Health}, vol.~6, no.~2, pp.~e142--e154, 2024.

\bibitem{taylor2024predicting}
T.~A.~Taylor, M.~RIFIER, and K.~L.~Evans, ``Predicting functional decline in amyotrophic lateral sclerosis using SHAP-enhanced gradient boosting,'' \textit{IEEE Journal of Biomedical and Health Informatics}, vol.~28, no.~4, pp.~2115--2126, 2024.

\bibitem{swoboda2005natural}
K.~G.~Swoboda, T.~W.~Prior, C.~B.~Scott, \textit{et al.}, ``Natural history of denervation in SMA: Relation to age, SMN2 copy number, and motor function,'' \textit{Annals of Neurology}, vol.~57, no.~5, pp.~704--712, 2005.

\bibitem{darras2021nusinersen}
B.~T.~Darras, M.~Masson, M.~P.~Mazurkiewicz, \textit{et al.}, ``Nusinersen in children and young adults with spinal muscular atrophy,'' \textit{Annals of Clinical and Translational Neurology}, vol.~8, no.~1, pp.~79--89, 2021.

\bibitem{barrera2022biomarkers}
J.~M.~Barrera, C.~A.~Vance, and E.~M.~Lombardi, ``Neurofilament light chain dynamics in infantile-onset spinal muscular atrophy,'' \textit{Neurology: Genetics}, vol.~8, no.~4, p.~e702, 2022.

\bibitem{alzaidy2023gene}
S.~Al-Zaidy, T.~P.~Mendell, and R.~Shelley, ``Long-term safety and durability of onasemnogene abeparvovec in spinal muscular atrophy,'' \textit{Lancet Neurology}, vol.~22, no.~11, pp.~1012--1023, 2023.

\bibitem{khatri2024multimodal}
I.~A.~Khatri, S.~U.~Hassan, and P.~R.~Narayanan, ``Predictive modeling in rare pediatric disorders: Challenges and opportunities for explainable machine learning,'' \textit{IEEE Reviews in Biomedical Engineering}, vol.~17, pp.~312--325, 2024.

\end{thebibliography}

\clearpage

% ==========================================
% TABLES SECTION: 8 COMPREHENSIVE IEEE TABLES
% ==========================================

% PAGE 17: Table 1 & Table 2
\begin{table*}[!t]
\caption{Baseline Demographic, Clinical, Genetic, and Motor Characteristics of the Cohort ($N = 350$)}
\label{tab:demographics}
\centering
\begin{tabular}{lcccc}
\toprule
\textbf{Variable / Clinical Parameter} & \textbf{Total Cohort ($N = 350$)} & \textbf{SMA Type 1 ($n = 154$)} & \textbf{SMA Type 2 ($n = 196$)} & \textbf{$p$-value} \\
\midrule
Sex (Male / Female), $n$ (\%) & 182 / 168 (52.0 / 48.0) & 81 / 73 (52.6 / 47.4) & 101 / 95 (51.5 / 48.5) & 0.842 \\
Gestational Age (weeks), mean $\pm$ SD & $39.1 \pm 1.4$ & $39.0 \pm 1.5$ & $39.2 \pm 1.3$ & 0.284 \\
Birth Weight (kg), mean $\pm$ SD & $3.24 \pm 0.45$ & $3.18 \pm 0.48$ & $3.29 \pm 0.42$ & 0.082 \\
Age at Symptom Onset (months), mean $\pm$ SD & $7.27 \pm 5.42$ & $2.42 \pm 1.28$ & $11.08 \pm 4.25$ & $< 0.001^{***}$ \\
Treatment Initiation Delay (months), mean $\pm$ SD & $3.81 \pm 2.45$ & $2.12 \pm 1.15$ & $5.14 \pm 2.38$ & $< 0.001^{***}$ \\
Chronological Age at Treatment Start (months), mean $\pm$ SD & $11.08 \pm 6.12$ & $4.54 \pm 1.82$ & $16.18 \pm 4.85$ & $< 0.001^{***}$ \\
\textit{SMN2} Copy Number, $n$ (\%) & & & & $< 0.001^{***}$ \\
\quad 2 Copies & 132 (37.7\%) & 132 (85.7\%) & 0 (0.0\%) & \\
\quad 3 Copies & 174 (49.7\%) & 22 (14.3\%) & 152 (77.6\%) & \\
\quad 4 Copies & 44 (12.6\%) & 0 (0.0\%) & 44 (22.4\%) & \\
Baseline CMAP Amplitude (mV), mean $\pm$ SD & $1.76 \pm 1.12$ & $0.82 \pm 0.48$ & $2.50 \pm 1.05$ & $< 0.001^{***}$ \\
Baseline Motor Score (Raw), mean $\pm$ SD & --- & $23.6 \pm 8.4$ (CHOP) & $28.4 \pm 9.6$ (HFMSE) & N/A \\
Normalized Baseline Motor Score (\%), mean $\pm$ SD & $39.9 \pm 13.8$ & $36.9 \pm 13.1$ & $43.0 \pm 14.5$ & $< 0.001^{***}$ \\
Nutritional Z-Score (BMI / WFA), mean $\pm$ SD & $-0.74 \pm 0.95$ & $-1.18 \pm 0.82$ & $-0.39 \pm 0.91$ & $< 0.001^{***}$ \\
Daily Baseline NIV Usage (hours/day), mean $\pm$ SD & $2.35 \pm 3.10$ & $3.55 \pm 3.82$ & $1.41 \pm 1.95$ & $< 0.001^{***}$ \\
Scoliosis Cobb Angle (degrees), mean $\pm$ SD & $11.8 \pm 8.4$ & $6.2 \pm 3.6$ & $16.2 \pm 8.5$ & $< 0.001^{***}$ \\
Treatment Regimen Modality, $n$ (\%) & & & & $< 0.001^{***}$ \\
\quad Nusinersen (Spinraza) & 158 (45.1\%) & 65 (42.2\%) & 93 (47.4\%) & \\
\quad Risdiplam (Evrysdi) & 117 (33.4\%) & 43 (27.9\%) & 74 (37.8\%) & \\
\quad Onasemnogene Abeparvovec (Zolgensma) & 57 (16.3\%) & 38 (24.7\%) & 19 (9.7\%) & \\
\quad Untreated Natural History Controls & 18 (5.1\%) & 8 (5.2\%) & 10 (5.1\%) & \\
\bottomrule
\end{tabular}
\end{table*}

\begin{table*}[!t]
\caption{Definition and Description of Machine Learning Multimodal Model Features}
\label{tab:features}
\centering
\begin{tabular}{lllll}
\toprule
\textbf{Feature Identifier} & \textbf{Clinical Domain} & \textbf{Type} & \textbf{Unit} & \textbf{Physiological Rationale and Relevance} \\
\midrule
\texttt{SMN2\_Copies} & Genomics & Discrete Integer & Count (2, 3, 4) & Primary phenotypic modifier governing basal full-length SMN expression. \\
\texttt{Age\_Onset\_mo} & Chronology & Continuous & Months & Reflects biological pace of neurodegenerative manifestation. \\
\texttt{Treatment\_Delay\_mo} & Chronology & Continuous & Months & Duration of unmitigated spinal motor neuron apoptosis prior to therapy. \\
\texttt{Age\_Treatment\_Start\_mo} & Chronology & Continuous & Months & Patient age at first therapeutic dose (\texttt{Age\_Onset} + \texttt{Treatment\_Delay}). \\
\texttt{Baseline\_CMAP\_mV} & Neurophysiology & Continuous & Millivolts (mV) & Electrophysiological proxy of viable motor units available for reinnervation. \\
\texttt{Normalized\_Base\_Score} & Motor Function & Continuous & Percent (\%) & Baseline motor score normalized to scale maximum (CHOP-INTEND or HFMSE). \\
\texttt{Nutritional\_ZScore} & Anthropometry & Continuous & Z-score & WHO age-adjusted weight/BMI Z-score; surrogate for bulbar integrity. \\
\texttt{NIV\_Hours\_Daily} & Respiratory & Continuous & Hours / day & Daily duration of non-invasive bi-level positive airway pressure support. \\
\texttt{Scoliosis\_Cobb\_deg} & Orthopedic & Continuous & Degrees ({\circ}$) & Major curve Cobb angle; reflects axial muscular weakness. \\
\texttt{SMA\_Type} & Phenotypic Subtype & Binary & Type 1 / Type 2 & Functional motor milestone categorization (sitters vs. non-sitters). \\
\texttt{Treatment} & Therapeutics & Categorical & 4 Classes & DMT agent (Nusinersen, Risdiplam, Onasemnogene, or Untreated Control). \\
\bottomrule
\end{tabular}
\end{table*}

\clearpage

% PAGE 18: Table 3 & Table 4
\begin{table*}[!t]
\caption{Longitudinal Motor Function Measurements Across Clinical Evaluation Intervals}
\label{tab:longitudinal}
\centering
\begin{tabular}{lccccccc}
\toprule
\textbf{Cohort Strata} & \textbf{Baseline} & \textbf{6 Months} & \textbf{12 Months} & \textbf{18 Months} & \textbf{24 Months} & \textbf{Mean $\Delta$ at 24m} & \textbf{MCID (\%)} \\
\midrule
\textbf{Full Cohort ($N=350$)} [Norm. \%] & $39.9 \pm 13.8$ & $44.1 \pm 14.2$ & $47.8 \pm 14.8$ & $49.5 \pm 15.0$ & $50.3 \pm 15.2$ & $+10.4 \pm 5.2\%$ & 72.6\% \\
\textbf{SMA Type 1 ($n=154$)} [CHOP-INTEND] & $23.6 \pm 8.4$ & $26.8 \pm 8.6$ & $29.4 \pm 9.1$ & $30.6 \pm 9.2$ & $31.2 \pm 9.4$ & $+7.6 \pm 4.8\text{ pts}$ & 76.6\% \\
\quad \textit{SMN2} = 2 Copies ($n=132$) & $22.2 \pm 7.9$ & $25.2 \pm 8.1$ & $27.8 \pm 8.6$ & $28.9 \pm 8.7$ & $29.5 \pm 8.9$ & $+7.3 \pm 4.6\text{ pts}$ & 75.0\% \\
\quad \textit{SMN2} = 3 Copies ($n=22$) & $31.8 \pm 6.2$ & $36.4 \pm 6.5$ & $39.0 \pm 6.8$ & $40.8 \pm 7.0$ & $41.4 \pm 7.2$ & $+9.6 \pm 5.2\text{ pts}$ & 86.4\% \\
\textbf{SMA Type 2 ($n=196$)} [HFMSE] & $28.4 \pm 9.6$ & $30.4 \pm 9.5$ & $31.8 \pm 9.8$ & $32.4 \pm 10.0$ & $32.8 \pm 10.1$ & $+4.4 \pm 3.2\text{ pts}$ & 69.4\% \\
\quad \textit{SMN2} = 3 Copies ($n=152$) & $26.6 \pm 9.4$ & $28.5 \pm 9.2$ & $29.8 \pm 9.5$ & $30.4 \pm 9.7$ & $30.8 \pm 9.8$ & $+4.2 \pm 3.1\text{ pts}$ & 67.8\% \\
\quad \textit{SMN2} = 4 Copies ($n=44$) & $34.6 \pm 7.1$ & $36.8 \pm 7.2$ & $38.8 \pm 7.5$ & $39.4 \pm 7.6$ & $39.8 \pm 7.8$ & $+5.2 \pm 3.4\text{ pts}$ & 75.0\% \\
\bottomrule
\end{tabular}
\end{table*}

\begin{table*}[!t]
\caption{Machine Learning Algorithms, Optimization Specifications, and Bayesian Hyperparameters}
\label{tab:hyperparams}
\centering
\begin{tabular}{llll}
\toprule
\textbf{Algorithm Architecture} & \textbf{Framework} & \textbf{Optimized Settings (Inner 5-Fold Bayesian Search)} & \textbf{Search Parameter Space Evaluated} \\
\midrule
\textbf{Gradient Boosting Regressor (GBM)} & Scikit-Learn & \texttt{n\_estimators=120}, \texttt{max\_depth=3}, \texttt{lr=0.08}, \texttt{subsample=0.85} & \texttt{n\_est} [50--300], \texttt{depth} [2--8], \texttt{lr} [0.01--0.2] \\
\textbf{Random Forest Regressor} & Scikit-Learn & \texttt{n\_estimators=150}, \texttt{max\_depth=6}, \texttt{min\_split=4}, \texttt{max\_feat='sqrt'} & \texttt{n\_est} [50--500], \texttt{depth} [3--15], \texttt{split} [2--10] \\
\textbf{Support Vector Regressor (SVR)} & Scikit-Learn & \texttt{kernel='rbf'}, \texttt{C=5.0}, \texttt{epsilon=0.20}, \texttt{gamma='scale'} & \texttt{C} [0.1--50], \texttt{epsilon} [0.01--1.0], \texttt{gamma} [auto, scale] \\
\textbf{Multi-Layer Perceptron (MLP)} & Scikit-Learn & \texttt{hidden=(64, 32)}, \texttt{activation='relu'}, \texttt{alpha=0.01}, \texttt{lr=0.001} & \texttt{layers} [(32,), (64, 32), (128, 64)], \texttt{alpha} [1e-4--0.1] \\
\textbf{ElasticNet Regression} & Scikit-Learn & \texttt{alpha=0.10}, \texttt{l1\_ratio=0.50}, \texttt{max\_iter=2000}, \texttt{tol=1e-4} & \texttt{alpha} [0.001--10], \texttt{l1\_ratio} [0.1--0.9] \\
\textbf{Ridge Regression} & Scikit-Learn & \texttt{alpha=1.00}, \texttt{solver='auto'}, \texttt{tol=1e-4} & \texttt{alpha} [0.01--100] \\
\bottomrule
\end{tabular}
\end{table*}

\clearpage

% PAGE 19: Table 5 & Table 6
\begin{table*}[!t]
\caption{Comprehensive Machine Learning Predictive Performance Comparison (10-Fold Repeated Nested Cross-Validation)}
\label{tab:comparison}
\centering
\begin{tabular}{lccccc}
\toprule
\textbf{Model Architecture} & \textbf{$R^2$ Score (Mean [95\% CI])} & \textbf{RMSE (Points)} & \textbf{MAE (Points)} & \textbf{Pearson $r$ ($p$-value)} & \textbf{Rank} \\
\midrule
\textbf{Gradient Boosting Regressor (GBM)} & \textbf{0.840 [0.816 -- 0.864]} & \textbf{1.77 $\pm$ 0.21} & \textbf{1.37 $\pm$ 0.16} & \textbf{0.926 ($< 0.001$)} & \textbf{1 (Best)} \\
Random Forest Regressor & 0.812 [0.784 -- 0.840] & 1.92 $\pm$ 0.24 & 1.51 $\pm$ 0.18 & 0.908 ($< 0.001$) & 2 \\
Support Vector Regressor (SVR, RBF) & 0.748 [0.716 -- 0.780] & 2.22 $\pm$ 0.26 & 1.76 $\pm$ 0.20 & 0.871 ($< 0.001$) & 3 \\
Multi-Layer Perceptron (MLP) & 0.722 [0.684 -- 0.760] & 2.34 $\pm$ 0.29 & 1.84 $\pm$ 0.22 & 0.854 ($< 0.001$) & 4 \\
ElasticNet Regression & 0.684 [0.648 -- 0.720] & 2.50 $\pm$ 0.31 & 1.98 $\pm$ 0.24 & 0.832 ($< 0.001$) & 5 \\
Ridge Regression & 0.672 [0.634 -- 0.710] & 2.55 $\pm$ 0.32 & 2.02 $\pm$ 0.25 & 0.824 ($< 0.001$) & 6 \\
\bottomrule
\end{tabular}
\end{table*}

\begin{table*}[!t]
\caption{SHAP-Based Global Feature Importance Ranking and Attribution Metrics}
\label{tab:shap}
\centering
\begin{tabular}{lclll}
\toprule
\textbf{Feature Description} & \textbf{Mean $|\text{SHAP}|$ ($E[|\phi|]$, pts)} & \textbf{Directional Impact Trend} & \textbf{Top Interaction} & \textbf{Mechanistic Pathophysiological Interpretation} \\
\midrule
\textbf{Treatment Initiation Delay (mo)} & \textbf{2.14} & Strongly Negative with Delay & Baseline CMAP & Irreversible spinal motor neuron loss prior to rescue. \\
\textbf{Baseline CMAP Amplitude (mV)} & \textbf{1.82} & Positive with Higher Amplitude & Treatment Delay & Functional motor unit pool capable of terminal reinnervation. \\
\textbf{Normalized Baseline Motor Score (\%)} & \textbf{1.45} & Non-Linear / Ceiling at High & \textit{SMN2} Copies & Defines functional headroom; exhibits psychometric ceiling. \\
\textbf{SMN2 Copy Number} & \textbf{1.28} & Positive with Higher Copy Count & Baseline Score & Primary genomic modifier governing endogenous SMN levels. \\
\textbf{Treatment: Onasemnogene Abeparvovec} & \textbf{1.12} & Strongly Positive Impact & Age at Onset & Rapid, systemic SMN replacement drives early motor recovery. \\
\textbf{Age at Symptom Onset (mo)} & \textbf{0.94} & Bimodal / Positive for Later & \textit{SMN2} Copies & Surrogate for biological aggressiveness and phenotype type. \\
\textbf{Daily Baseline NIV Usage (hrs)} & \textbf{0.58} & Negative with Higher Hours & Baseline Score & Marker of respiratory diaphragmatic failure and debility. \\
\textbf{Scoliosis Cobb Angle ({\circ}$)} & \textbf{0.42} & Negative with Severe Angle & Age at Onset & Biomechanical restriction impairing axial motor tasks. \\
\textbf{Nutritional Z-Score} & \textbf{0.34} & Negative with Severe Underweight & Daily NIV Hours & Marker of bulbar feeding failure and systemic cachexia. \\
\textbf{Treatment: Untreated Control} & \textbf{0.31} & Extremely Negative ($-6.60$ pts) & All Variables & Natural history trajectory of unmitigated neurodegeneration. \\
\bottomrule
\end{tabular}
\end{table*}

\clearpage

% PAGE 20: Table 7 & Table 8
\begin{table*}[!t]
\caption{SMA Type 1 and Type 2 Subgroup Model Performance Stratification}
\label{tab:subgroup}
\centering
\begin{tabular}{lcccccc}
\toprule
\textbf{Clinical / Genetic Subgroup} & \textbf{Patient Count ($n$)} & \textbf{$R^2$ Score} & \textbf{RMSE (Points)} & \textbf{MAE (Points)} & \textbf{Pearson $r$} & \textbf{Concordance $\rho_c$} \\
\midrule
\textbf{Full Cohort} & \textbf{350} & \textbf{0.840} & \textbf{1.77} & \textbf{1.37} & \textbf{0.926} & \textbf{0.918} \\
\textbf{SMA Type 1 Subgroup} & 154 & 0.852 & 1.82 & 1.41 & 0.931 & 0.924 \\
\quad \textit{SMN2} = 2 Copies & 132 & 0.848 & 1.80 & 1.39 & 0.928 & 0.920 \\
\quad \textit{SMN2} = 3 Copies & 22 & 0.868 & 1.94 & 1.52 & 0.942 & 0.935 \\
\textbf{SMA Type 2 Subgroup} & 196 & 0.828 & 1.72 & 1.33 & 0.916 & 0.908 \\
\quad \textit{SMN2} = 3 Copies & 152 & 0.835 & 1.74 & 1.35 & 0.920 & 0.912 \\
\quad \textit{SMN2} = 4 Copies & 44 & 0.810 & 1.66 & 1.28 & 0.905 & 0.898 \\
\bottomrule
\end{tabular}
\end{table*}

\begin{table*}[!t]
\caption{Sensitivity and Robustness Evaluation Across Methodological Perturbations}
\label{tab:sensitivity}
\centering
\begin{tabular}{lcccc}
\toprule
\textbf{Experimental Perturbation Condition} & \textbf{$R^2$ (Mean $\pm$ SD)} & \textbf{RMSE (Points)} & \textbf{MAE (Points)} & \textbf{Performance Delta ($\Delta R^2$)} \\
\midrule
\textbf{Baseline Primary Model (10-Fold CV)} & \textbf{0.840 $\pm$ 0.038} & \textbf{1.77 $\pm$ 0.21} & \textbf{1.37 $\pm$ 0.16} & \textbf{Reference} \\
Exclusion of CMAP Amplitude & $0.782 \pm 0.045$ & $2.06 \pm 0.25$ & $1.62 \pm 0.19$ & $-0.058$ \\
Exclusion of Treatment Delay & $0.765 \pm 0.048$ & $2.14 \pm 0.27$ & $1.69 \pm 0.21$ & $-0.075$ \\
Exclusion of \textit{SMN2} Copy Number & $0.804 \pm 0.041$ & $1.96 \pm 0.23$ & $1.54 \pm 0.18$ & $-0.036$ \\
Addition of 15\% Gaussian Feature Noise & $0.818 \pm 0.042$ & $1.89 \pm 0.23$ & $1.48 \pm 0.18$ & $-0.022$ \\
50 Random Seed Monte Carlo Iterations & $0.838 \pm 0.012$ & $1.78 \pm 0.11$ & $1.38 \pm 0.09$ & $-0.002$ \\
\bottomrule
\end{tabular}
\end{table*}

\end{document}
